\chapter{Полярная мера, BV-фактор и допустимое окно барьера потери ранга}

Предыдущая проверка точно локализовал проблему: плоский интеграл по мосту
имеет на чистой страте остаточную расходимость $dt/t$. Теперь проверяется,
может ли она исчезнуть после правильного полярного разложения, удаления
ортогональных орбит или введения естественной индуцированной меры на
семейных состояниях.

\section{Полярное разложение не меняет интеграл}

Для общего поперечного блока $C\in M_2(\mathbb R)$ запишем сингулярное
разложение
\begin{equation}
 C=U\operatorname{diag}(s_1,s_2)V^T,
 \qquad U,V\in O(2),\qquad s_i\geq0.
\end{equation}
Плоская мера имеет вид
\begin{equation}
 dC=\mathrm{const}\,
 |s_1^2-s_2^2|\,ds_1ds_2\,d\mu(U)d\mu(V).
 \label{eq:v6-real-m2-svd-measure}
\end{equation}
Это стандартный вещественный матричный якобиан; см.
\path{arXiv:1701.04505}. При $s_i=t a_i$ радиальная часть снова
масштабируется как $t^3dt$. После поперечного гауссова множителя $t^{-4}$
остаётся прежнее $dt/t$.

Если $O(2)\times O(2)$ считать калибровочной орбитой, формула Фаддеева--
Попова оставляет якобиан $|s_1^2-s_2^2|$ на сечении. Деление на конечный
объём компактной группы удаляет только постоянный множитель. Оно не
превращает $t^3dt$ в двумерную плоскую меру сингулярных чисел.

\begin{theorem}[Полярный/BV-якобиан не насыщает долину]
Полярное разложение, сингулярные координаты и стандартное удаление
компактных ортогональных орбит являются точной переписью исходной меры.
Они сохраняют остаточную асимптотику $dt/t$ и не создают зависящего от
состояния барьера $-\nu\log\det R$.
\end{theorem}

Контрактные BV-пары также не могут изменить классическое множество
минимумов: они организуют калибровочный фактор и квантовые детерминанты,
но новый классический потенциал потребовал бы новых физических данных.

\section{Естественный кандидат из индуцированной меры состояний}

Рассмотрим отдельно стандартную вещественную индуцированную меру. Если
$X\in M_{3\times K}(\mathbb R)$ имеет плоскую вращательно-инвариантную
меру и
\begin{equation}
 R=\frac{XX^T}{\Tr(XX^T)},
\end{equation}
то относительно плоской меры на вещественных матрицах плотности возникает
фактор
\begin{equation}
 d\mu_K(R)\propto
 (\det R)^{(K-4)/2}\,dR,
 \qquad R>0,\qquad\Tr R=1.
 \label{eq:v6-real-induced-state-measure}
\end{equation}
В эффективном действии это соответствует
\begin{equation}
 \mathcal F_{\rm meas}(R)=-\nu_K\log\det R,
 \qquad \nu_K=\frac{K-4}{2}.
 \label{eq:v6-induced-logdet-barrier}
\end{equation}
Это Уишарта-якобиан, а не свободно назначенный потенциал.

Однако текущий родитель не содержит матрицу очищения $3\times K$. Особенно
важно, что встречающееся в Томе V число пять является виртуальной разностью
$20-15$. Том V отдельно доказал отсутствие канонического положительного
проектора ранга пять в двадцатимерном углу. Поэтому выбор $K=5$ пока не
является выведенной микроскопической мерой.

\section{Совместимость с самозапуском}

Даже если временно проверить общий член
$-\nu\log\det R$, он не решает задачу. Около
$R=I_3/3+\delta R$, $\Tr\delta R=0$, получаем
\begin{equation}
 -\nu\log\det R
 =\mathrm{const}+\frac{9\nu}{2}\Tr(\delta R)^2+O(\delta R^3).
\end{equation}
До этого члена полный локальный коэффициент был $-51/112$. Следовательно,
самозапуск сохраняется только при
\begin{equation}
 -\frac{51}{112}+\frac{9\nu}{2}<0
 \quad\Longleftrightarrow\quad
 \nu<\frac{17}{168}.
 \label{eq:v6-logdet-instability-upper-bound}
\end{equation}

На чистой границе нельзя продолжать однопетлевую формулу: именно там её
массы обращаются в нуль. Нужно использовать непертурбативный степенной
счёт предыдущей проверки. Для
$R_\varepsilon=\operatorname{diag}(1-2\varepsilon,
\varepsilon,\varepsilon)$ квартичная часть обрезает плоскую долину при
$t\sim\varepsilon^{-1/4}$, поэтому
\begin{equation}
 Z_{\rm br}(R_\varepsilon)
 \sim A\int^\infty\frac{dt}{t}e^{-c\varepsilon t^4}
 =\frac A4 E_1(c\varepsilon)
 \sim\frac A4\log\frac1\varepsilon.
 \label{eq:v6-nonperturbative-boundary-log}
\end{equation}
Следовательно,
$\mathcal F_{\rm br}\sim-\log\log(1/\varepsilon)$, тогда как
$-\nu\log\det R_\varepsilon\sim2\nu\log(1/\varepsilon)$.
Любое $\nu>0$ асимптотически подавляет границу.

\begin{theorem}[Допустимое окно слабого барьера]
В непертурбативном степенном приближении барьер
$-\nu\log\det R$ совместим с локальным самозапуском ровно в окне
\begin{equation}
 \boxed{0<\nu<\frac{17}{168}}.
 \label{eq:v6-admissible-fractional-barrier-window}
\end{equation}
Полярная мера даёт $\nu=0$ и не насыщает переход. Стандартная вещественная
индуцированная мера с целым $K$ также не попадает в окно: ближайший
положительный коэффициент при $K=5$ равен $1/2$ и уже уничтожает
отрицательную моду.
\end{theorem}

Тем самым закрыт не любой логарифмический барьер, а его происхождение из
стандартной полярной или уишартовской меры. Остаётся конкретный вопрос:
может ли относительный детерминант или нормированный проектный след дать
дробный коэффициент внутри допустимого окна
\eqref{eq:v6-admissible-fractional-barrier-window}.
Числа $1/7$ и $1/14$ нельзя подставлять по сходству; их возможное появление
должно следовать из полного комплекса мод.

\section{Физический смысл для картины кристаллизации}

Обычная кинематическая мера не завершает кристаллизацию. Но
непертурбативная проверка не закрывает детерминантный маршрут полностью:
она выделяет узкий диапазон слабого дробного барьера. Такой коэффициент
должен возникнуть из настоящего относительного комплекса мод, а не из
произвольной меры.

Это уточняет раннюю идею проекта. Материя может читаться как устойчивые
дефекты кристаллизованной операторной среды, но сначала теория обязана
вывести конечную упорядоченную фазу. Само слово ``кристалл'' не заменяет
этот вариационный шаг.

\begin{nogocriterion}[Запрет назначения дробного барьера]
Нельзя вводить $-\nu\log\det R$, выбирать размерность очищения $K$ или
удалять орбитальный якобиан только ради получения конечного минимума.
Следующая проверка обязана вывести коэффициент $0<\nu<17/168$ из уже имеющегося
относительного детерминанта/BV-комплекса либо закрыть этот последний мерный
маршрут.
\end{nogocriterion}

\begin{protocol}
\begin{enumerate}
 \item полярный якобиан $M_2(\mathbb R)$:
 \textbf{$|s_1^2-s_2^2|$};
 \item радиальная степень после SVD: \textbf{$t^3dt$};
 \item остаток после поперечного интеграла: \textbf{$dt/t$};
 \item стандартный FP/BV-фактор как лекарство: \textbf{нет};
 \item индуцированный Уишарта-барьер:
 \textbf{$\nu_K=(K-4)/2$};
 \item каноническое очищение $3\times K$ в текущем родителе:
 \textbf{не выведено};
 \item условие сохранения самозапуска:
 \textbf{$\nu<17/168$};
 \item непертурбативный рост $Z_{\rm br}$ у чистой вершины:
 \textbf{$\log(1/\varepsilon)$};
 \item условие подавления чистой вершины:
 \textbf{$\nu>0$};
 \item допустимое окно слабого барьера:
 \textbf{$0<\nu<17/168$};
 \item стандартное целочисленное Уишарта-происхождение:
 \textbf{не попадает в окно};
 \item следующая проверка:
 \textbf{происхождение дробного детерминантного коэффициента}.
\end{enumerate}
\end{protocol}

Машинный сертификат сохранён в
\path{s2t_v6_polar_bv_rank_loss_barrier_gate_results.json}.